% day1_seg4_prob/sample_space.tex
\begin{frame}
  \onslide<1->{
	\frametitle{Probability} 
	\framesubtitle{Sample Spaces}
}
\begin{itemize}
  \onslide<2->{
\item  Experiments 
}
  \onslide<3->{
\item Sample Space:  $\Omega$  
\begin{itemize}
 \item consists of all possible outcomes $\omega$.
 \end{itemize}
}
  \onslide<4->{
\item Events:  Subsets of $\Omega$ forming a $\sigma$-algebra.  Not necessarily \textit{all} subsets.  
}
\end{itemize}
\end{frame}

\begin{frame}
  \onslide<1->{
	\frametitle{Probability} 
	\framesubtitle{Examples:  Experiment and Sample Space}
}
\begin{itemize}
  \onslide<2->{
\item 1.  Experiment:  Single coin toss
  \begin{itemize}
  \item Sample Space:  $\Omega=\set{H,T}$
  \end{itemize}  
}
  \onslide<3->{
\item 2. Experiment:  Toss coin twice
  \begin{itemize}
  \item Sample Space:  $\Omega=\set{HH,HT,TH,TT}$
  \end{itemize}
}
\end{itemize}
\end{frame}

\begin{frame}
  \onslide<1->{
	\frametitle{Probability} 
	\framesubtitle{$\sigma$-Algebras}
}
\begin{itemize}
  \onslide<2->{
\item In probability, we are often less concerned with individual outcomes and more with sets of outcomes.  
}
  \onslide<3->{
\item Sets of outcomes, i.e., subsets of the sample space $\Omega$, are called "events".    
}
  \onslide<4->{
\item Not every subset of the sample space $\Omega$ qualifies as an event.     
}
  \onslide<5->{
\item Qualified events form an "event algebra", known in mathematics as a \textbf{$\sigma$-algebra}  
}
\end{itemize}
\end{frame}

\begin{frame}
  \onslide<1->{
	\frametitle{Probability} 
	\framesubtitle{$\sigma$-Algebras}
}
\begin{itemize}
  \onslide<2->{
\item  A $\sigma$-algebra $\mathcal{M}$ is a
 collection of subsets of $\Omega$ that is closed under arbitrary unions and
 intersections, set complementation, and which also contains the whole
 space $\Omega$ as well. 
}
  \onslide<3->{
\item A $\sigma$-algebra can be expressed in structural notation as\\
 $<\mathcal{M},\cup,\cap,\:',\emptyset,\Omega>$  
}
  \onslide<4->{
\item $\Omega$ is the certain event; $\emptyset$ is the impossible event
}
\end{itemize}
\end{frame}

\begin{frame}
  \onslide<1->{
	\frametitle{Probability} 
	\framesubtitle{Examples: $\sigma$-algebras}
}
\begin{itemize}
  \onslide<2->{
\item Double coin toss:
}
  \onslide<3->{
  \begin{itemize}
  \item Sample Space: $\Omega=\set{HH,HT,TH,TT}$
  \item $\sigma$-algebra: $\mathcal{M}=\set{\emptyset, \set{HH},\set{HT,TH,TT},\Omega}$
   \item Note that for this $\sigma$-algebra, the set $\set{TT}$ is not considered an event
  \end{itemize}  
}  
   \onslide<4->{
  \item The power set $\mathcal{P}(\Omega)$, consisting of all subsets of $\Omega$, is a $\sigma$-algebra
}
   \onslide<5->{
  \item The set consisting of just the empty set and the sample space itself is also a $\sigma$-algebra:  
  $\mathcal{M}=\set{\emptyset,\Omega}$
}
\end{itemize}
\end{frame}

\begin{frame}
  \onslide<1->{
	\frametitle{Probability} 
	\framesubtitle{Mutually Exclusive Events}
}
\begin{itemize}
  \onslide<2->{
\item Events $A,B\in\mathcal{M}$ are called \textbf{\textit{mutually exclusive}}, 
if they are disjoint as sets, i.e., if $A\cap B=\emptyset$
}
  \onslide<3->{
\item In previous example, events $\set{HH}$ and $\set{HT,TH,TT}$ are mutually exclusive
}
\end{itemize}
\end{frame}

